% Generated draft: requires mathematical and source review.
\documentclass[11pt]{article}
\usepackage[margin=25mm]{geometry}
\usepackage{fontspec}
\setmainfont{Latin Modern Roman}
\usepackage{amsmath,amssymb,mathtools,unicode-math}
\setmathfont{Latin Modern Math}
\usepackage{enumitem,needspace,xurl,hyperref}
\setlength{\parindent}{0pt}
\setlength{\parskip}{5pt}
\setlength{\emergencystretch}{3em}
\newcommand{\N}{\mathbb{N}}
\newcommand{\Z}{\mathbb{Z}}
\newcommand{\Q}{\mathbb{Q}}
\newcommand{\R}{\mathbb{R}}
\newcommand{\C}{\mathbb{C}}
\newcommand{\F}{\mathbb{F}}
\newcommand{\A}{\mathbb{A}}
\newcommand{\PP}{\mathbb{P}}
\newcommand{\E}{\mathbb{E}}
\newcommand{\eps}{\varepsilon}
\newcommand{\dd}{\,\mathrm{d}}
\begin{document}
\section*{UnsolvedMath ID: 800056}
\textbf{Source problem code:} catalogue entry
\subsection{Definitions and mathematical statement}
Let \(M\) be a smooth manifold with an almost-complex structure \(J\colon TM\to TM\), satisfying \(J^2=-\mathrm{id}_{TM}\). A compatible Hermitian structure is represented by a Riemannian metric \(g\) satisfying
\[
g(JX,JY)=g(X,Y)
\]
for all tangent vectors \(X,Y\). Write \(\dim_{\R}M=2m\); the source writes \(M^n\) without specifying its dimension convention.

The source refers to numbers \(h^{p,q}\) ``defined as above'', but the supplied record does not contain that preceding definition. Denote these unspecified numbers provisionally by \(h^{p,q}_g(M,J)\), for the bidegrees to which the missing definition applies.

The first question is whether, for every almost-complex manifold \((M,J)\) in the intended class and every pair of compatible Hermitian metrics \(g_0,g_1\), one has
\[
h^{p,q}_{g_0}(M,J)=h^{p,q}_{g_1}(M,J)
\]
for every relevant \((p,q)\).

If this independence fails, the source asks for another definition of numbers \(\widetilde h^{p,q}(M,J)\) determined solely by the almost-complex structure and generalizing the \(h^{p,q}\) of complex manifolds. Under the usual Dolbeault interpretation of complex Hodge numbers, the required agreement for integrable \(J\) would be
\[
\widetilde h^{p,q}(M,J)=\dim_{\C}H^{p,q}_{\bar\partial}(M,J),
\]
where
\[
H^{p,q}_{\bar\partial}(M,J)=
\frac{\ker\bigl(\bar\partial\colon\Omega^{p,q}(M)\to\Omega^{p,q+1}(M)\bigr)}
{\operatorname{im}\bigl(\bar\partial\colon\Omega^{p,q-1}(M)\to\Omega^{p,q}(M)\bigr)}
\]
and \(\Omega^{p,q}(M)\) denotes the space of smooth complex differential forms of type \((p,q)\). This interpretation of the extension requirement is conditional: the omitted preceding source text must be checked.

The supplied record does not specify compactness, boundary conditions, or the construction of \(h^{p,q}_g\). Consequently, it does not determine a complete precise formulation of the first question or establish that the displayed dimensions are finite.
\subsection{Short English statement}
For a fixed almost-complex structure, ask whether the source's numbers \(h^{p,q}\) are independent of the compatible Hermitian metric. If they are not, seek replacements determined by the almost-complex structure that recover the complex-manifold numbers. The defining passage and geometric hypotheses are missing, so this is a conditional statement-only draft requiring review.
\subsection{Sources}
\begin{itemize}
\item Local UnsolvedMath registry, supplied complete record with permanent ID \(800056\), exact source problem code \texttt{catalogue entry}, title ``Hirzebruch's 1954 Problem List --- Problem 5''. The registry attributes the entry to Kotschick (2013) and labels it open; that label is supplied metadata, not an independently verified mathematical assessment.
\item Kotschick, \emph{Updates on Hirzebruch's 1954 Problem List} (2013), Problem 5. Supplied source URL: \url{https://arxiv.org/abs/1305.4623}. The background records extraction by \texttt{arxiv\_latex} and an access date of 2026-10-16. These are supplied provenance details; no source visit or independent verification is asserted here.
\end{itemize}
\end{document}
